\documentclass[aps,prd,reprint,nofootinbib]{revtex4-2}
\usepackage{amsmath,amssymb,booktabs}
\usepackage[colorlinks=true,citecolor=blue,urlcolor=blue,linkcolor=blue]{hyperref}
% GENERATED by tools/evidence.py; source hashes in results.json
\newcommand{\result}[1]{\ifcsname result#1\endcsname\csname result#1\endcsname\else\PackageError{results}{Undefined result #1}{}\fi}
\expandafter\def\csname resultIntegralError\endcsname{0.00000509}
\expandafter\def\csname resultIntegralEstimate\endcsname{0.33332825}
\expandafter\def\csname resultIntegralReference\endcsname{0.33333333}
\expandafter\def\csname resultIntegralTolerance\endcsname{0.00001000}

\begin{document}
\title{Transparent agentic science: Documentation, reliability, and scientific judgment}
\author{R.~O'Shaughnessy}
\date{October 23, 2026}
\begin{abstract}
Artificial intelligence agents can assist with computation, debugging, literature searches, and manuscript preparation, but their usefulness depends on how their contributions are checked and interpreted. We describe a living guide to reproducible scientific work developed through human--agent collaboration in gravitational-wave astrophysics. The guide addresses four concerns: traceable numerical documentation, reliable calculations, accurate attribution to the literature, and readable scientific arguments. Generated result macros connect numerical statements to recorded computations; independent comparisons and adversarial review challenge the results; citation tools connect attributions to inspected sources. Human oversight remains essential throughout, particularly when physical intuition and plots reveal inconsistencies that automated checks miss. Writing presents the largest continuing difficulty: fluent and abundant prose can obscure priorities, context, and the central scientific contribution. Strong models, sustained revision, and explicit human direction are needed to improve the argument itself. The examples concern the main papers in two research repositories. The accompanying public repository is assembled entirely with AI from prior work with a human, and all its commits are performed by AI agents.
\end{abstract}
\maketitle

\section{Introduction}\label{sec:introduction}
The use and capabilities of artificial intelligence (AI) in science are expanding rapidly. Studies of the research literature document accelerating adoption across scientific fields~\cite{Hao2026}, while demonstrations of language-model systems connected to software and laboratory tools show how their role can extend from text generation to executing research tasks~\cite{Boiko2023}. In computational physics, an agent can translate a scientific question into code, run an analysis, inspect its output, and propose a revision within one working session. These capabilities make it possible to explore more cases and perform more cross-checks, but they also increase the amount of material a scientist must assess. The central challenge is to make the additional work scientifically useful and independently inspectable.

Journal and preprint policies identify several concerns that bear directly on this challenge. The Physical Review policy permits substantive AI assistance with disclosure, while retaining author responsibility for accuracy and verification~\cite{APS2026}. The guidance for the Astrophysical Journal (ApJ) and other American Astronomical Society (AAS) journals emphasizes supervision, disclosure, original work, and checks on data, figures, and references; it specifically warns that generated reference lists can be incomplete or fictitious~\cite{AAS2026}. arXiv likewise requires significant use of generative language tools to be reported and holds each author responsible for the content~\cite{arXivAI}. These concerns address concrete failure modes of poor or unchecked use: invented evidence, unreliable calculations, misleading attribution, and plausible prose that has not received scientific scrutiny. Disclosure makes the process visible, but a disclosed error remains an error.

In this paper, we present a living guide to conducting science reliably with agents. The guide translates these concerns into practices for documenting results, validating calculations, consulting the literature, and constructing a scientific argument. It draws on human--agent work on gravitational-wave population inference, illustrated by the main synthetic-universe and hierarchical-modeling papers~\cite{SyntheticUniverses,HierarchicalMethods}, together with a collection of reusable research skills. Our aim is to preserve the standards expected of any scientific investigation while using agents to make the supporting work easier to perform and inspect. The guide will change as tools improve and experience exposes new failure modes; its current recommendations are grounded in this collaboration rather than a controlled comparison of agents and human researchers.

A useful division of responsibility emerges from this experience. Computations, source retrieval, and consistency checks can often be delegated to tools with well-defined inputs and outputs. Choosing the scientific question, recognizing an implausible result, interpreting a limitation, and deciding which argument matters require scientific judgment. Agents can contribute to those decisions, including by challenging the investigator's expectations, but agreement among generated accounts cannot replace evidence. The distinction is especially important in writing, where a technically accurate draft may still fail to communicate the purpose or significance of the work.

This paper is organized as follows. Section~\ref{sec:documentation} discusses numerical documentation and provenance. Section~\ref{sec:reliability} considers reliable computation, adversarial review, and human oversight. Section~\ref{sec:literature} describes literature retrieval and citation checks. Section~\ref{sec:writing} addresses scientific writing, the component that currently requires the most sustained judgment. Section~\ref{sec:conclusions} summarizes the approach. Appendix~\ref{app:tooling} provides a small executable illustration of numerical consistency, and Appendix~\ref{app:resources} identifies the accompanying resources.

\section{Documenting numerical evidence}\label{sec:documentation}
A quantitative statement must refer to an identifiable calculation. A language model's ability to produce a plausible number is irrelevant to whether that number was measured or computed. The analysis must produce a saved result, with sufficient information to recover its inputs, method, uncertainty, and domain of validity. We organize this documentation around the connection from data to computation to scientific claim. This approach is modeled on the provenance and review practices used in LVK internal review, adapted here to research conducted with agents. The public LVK data-release documentation provides a related example of connecting reported quantities to identifiable data products~\cite{GWTCRelease}.

The essential mechanism is a generated numerical macro. The analysis writes a result record containing the raw value, units, uncertainty definition, input and configuration identities, and producing code. A deterministic tool converts this record into the quantities displayed in the manuscript. The abstract, text, table, and caption then refer to the same named result. A change to that result propagates to every use when the document is rebuilt. The same record can supply table cells and numerical annotations in plots, eliminating separate transcriptions that otherwise become inconsistent during revision.

Provenance must describe the calculation that actually ran. A source comment, directory name, or remembered command may describe an intended configuration rather than the effective configuration. Relevant settings therefore belong in the output itself, alongside the input identities, code revision, environment, seed where applicable, and termination status. Cached products require the same care: a calibration or selection change can invalidate a cache without changing its filename. Their identity must include the dependencies that determine their contents. Optimization results also require the initialization, restart schedule, checkpoints, and rule used to select the reported solution.

Documentation has a scientific purpose beyond bookkeeping. A historical comparison may intentionally retain its original sample, whereas a current inference must use the current analysis inputs. Both can be valid, provided that their domains are stated. A stale result, an explicitly frozen historical result, and a curated input constant are therefore distinct. Treating all older numbers as interchangeable invites errors: automatically updating a historical comparison can be as misleading as retaining a superseded measurement.

Tooling checks the consistency of this chain. A clean build regenerates numerical products and compares them with the recorded artifacts, identifies undefined result names, and flags computed quantities that bypass the result record. It must also fail when its required tools are absent. These checks prevent specified forms of transcription drift; they do not establish that the estimator or scientific interpretation is correct. Selecting the wrong result key remains a scientific error even if that key is used consistently throughout the paper. Appendix~\ref{app:tooling} illustrates both the utility and the limited guarantee of the mechanism.

The manuscript and its supporting notes serve different readers. Scientific notes retain derivations, rejected approaches, uncertainty studies, and the evidence behind changes of interpretation. The paper selects the material needed to understand and assess its claims. Keeping that distinction allows documentation to remain detailed without forcing the reader to reconstruct the argument from a chronological account of runs. The research record can be extensive while the manuscript remains focused.

\section{Establishing reliable results}\label{sec:reliability}
\subsection{Tools, comparisons, and uncertainty}
Agents are useful investigators when a question can be turned into a discriminating calculation. They can enumerate boundary cases, compare implementations, trace a discrepancy through several layers of code, and construct a regression test after locating a defect. In our experience, this is among their strongest contributions to scientific work. The requirement is that they use the relevant tools and inspect the resulting evidence. A numerical answer explained fluently by a model is not a substitute for executing a calculation, and a proposed test is not evidence that the test passed.

Validation is most informative when its alternatives fail for different reasons. Analytic answers and planted-truth examples check known outcomes. Limiting cases test physical conventions and normalization. Refinement studies expose approximation error, and independent seeds characterize stochastic variation. A second implementation can identify mistakes in the first, provided that their shared dependencies and assumptions are understood. The acceptance criterion should be chosen from the scientific question and its error budget, rather than adjusted after seeing the comparison.

For inference, reliable execution must be separated from reliable interpretation. A sampler can mix well while the parameter of interest remains unidentified or pressed against a prior boundary. A global likelihood improvement can coexist with failure in the subpopulation that motivates the claim. The main hierarchical-modeling paper illustrates both issues: it distinguishes a boundary-limited result from a measurement and uses a detected-population comparison to restrict the interpretation of a fitted model~\cite{HierarchicalMethods}. The relevant test follows the claim into the space where it has an observable consequence.

Finite numerical support deserves particular attention. Scatter across seeds does not by itself establish convergence with sample size, and more optimizer iterations cannot repair a simulation library that does not populate the relevant region. The main synthetic-universe paper accordingly treats an additional population component as a hypothesis when its available numerical support does not permit a reliable evidence comparison~\cite{SyntheticUniverses}. Recognizing such a limitation can be more useful than constructing another estimator around the same insufficient inputs.

\subsection{Adversarial review and human oversight}
Adversarial review asks how a result could fail. We distinguish a science review of the claim, assumptions, and domain from a code and numerics review of the implementation and its checks. Both examine a fixed candidate and a stated scientific specification. They seek counterexamples, convention errors, confounded comparisons, and unjustified extrapolations, with enough information to reproduce the objection. A reviewer should know the question and assumptions without being primed by the implementer's preferred explanation of why the answer must be correct.

The independence of a check depends on its evidence as well as its reviewer. Separate agent contexts can reduce agreement driven by a shared conversation, but they do not remove a common library error, input defect, or modeling assumption. The record should state which route was independent and for what purpose. A case used to repair or select the candidate becomes development evidence; a subsequent confirmation needs fresh cases where feasible. Deliberately incorrect inputs or isolated mutations can show that a check is capable of rejecting a relevant failure.

Human oversight is especially valuable when physical intuition identifies something that the existing tests did not ask about. A plot can make an implausible boundary behavior, displaced population, missing tail, or incorrect normalization immediately apparent. Such a judgment should initiate a test: identify the expected behavior, compare it with the plotted quantity, and trace any discrepancy through inputs, transformations, selection effects, and visualization code. The scientist's intuition is a source of hypotheses, not a desired answer that the agent must reproduce. A useful agent can also demonstrate when that intuition was wrong.

Plots therefore need direct scientific inspection, not merely confirmation that an image file exists. Axes, units, cuts, normalization, uncertainty bands, and the distinction between a likelihood and a probability density can all change the apparent conclusion. A summary statistic may conceal a shape discrepancy that is obvious in an overlay or residual plot. Visualization code must read the recorded data and preserve its meaning; an image generated to look like a scientific plot supplies no evidence. The AAS guidance makes the same practical distinction between supervised visualization of supplied data and invented scientific graphics~\cite{AAS2026}.

Review also governs revision. Findings remain recorded with their dispositions, including unresolved disagreement and objections rejected with evidence. When a claim fails, its supporting note is marked accordingly, the defeating test becomes a standing diagnostic, and every dependent statement is checked. The abstract, captions, conclusions, and related main paper can otherwise preserve an argument that the calculation no longer supports. This sweep is repeated after the final human or collaborator edit. Human involvement does not make a revision immune to staleness.

\section{Finding and citing relevant literature}\label{sec:literature}
Scientific attribution requires both a real source and a correct account of what it establishes. A model can invent a bibliographic entry, merge details from different papers, or attach a genuine citation to a claim that the cited work does not support. It can also supply a superficially credible but incomplete account of a field. The solution begins with the same tools and habits used by careful human researchers: search the literature, retrieve the source, inspect the relevant passage, and preserve its bibliographic identity.

Agents should use citation and discovery services such as ADS, arXiv, DOI resolution, and journal records. A reference manager can import verified metadata, while simple checks identify duplicate entries, undefined citation keys, missing identifiers, or mismatches in author, title, and year. These checks establish identity. Scientific support requires reading the cited equation, result, or discussion in context, including assumptions and limitations. A generated summary is a starting point for inspection rather than evidence that the source was read.

Relevance and coverage require additional judgment. Searches should follow the scientific ideas, methods, and competing interpretations rather than only the vocabulary suggested by the current draft. Citation chains and domain expertise help recover work that a generated list misses. The ApJ/AAS guidance explicitly recommends conventional literature tools, including ADS, alongside checks of whether suggested references exist and are relevant~\cite{AAS2026}. These are ordinary scholarly responsibilities; using an agent makes them easier to automate in part without making them optional.

The same standard applies to this guide. Journal requirements are cited to the current official policies, and the scientific examples refer to the main papers in their respective repositories. Repository links identify the supporting research context but currently require access; the preprints provide public scientific sources. We do not infer a claim's validity from the existence of either a repository or a preprint. The relevant evidence must still support the sentence for which it is cited.

\section{Writing a scientific argument}\label{sec:writing}
Writing is currently the part of this process that requires the most judgment. Agents can generate large amounts of grammatical prose quickly, yet that prose may give the wrong material the greatest weight. A draft can dwell on implementation minutiae, repeat caveats, miss the context that motivates the work, or describe a sequence of calculations without identifying the scientific contribution. These are failures of argument and priority. Correcting the wording alone leaves the underlying problem intact.

The human author must supply an explicit scientific brief: the question, audience, principal contribution, important comparisons, limits of the evidence, and intended thrust. This brief determines what belongs in the paper and what belongs in notes or supporting documentation. The agent can then help expose gaps, propose a structure, challenge a transition, or draft a bounded passage. The direction of work is from the scientific argument to the prose. A collection of polished paragraphs does not become a paper merely by placing headings between them.

Our practical recommendation is to use the strongest available models that have demonstrated useful performance on the task and to reserve substantial time for revision. Model quality matters for sustained reasoning about context and the relationship between sections; it does not eliminate the need to supply that context. The author may need to explain repeatedly why a particular result matters, why another belongs in an appendix, or why the apparent novelty is already established in the literature. These recommendations reflect experience in this collaboration, rather than a measured ranking of models or a guarantee about future systems.

Revision should first ask whether the argument is worth reading. The introduction must connect the broader scientific problem to the contribution. Methods must explain the assumptions and formalism needed to assess the result. Results must distinguish what was learned from what was merely run. Discussion must address implications and limitations at the appropriate scale. The abstract and conclusions should state the same central finding with a strength supported by the evidence. This is where agents often need the most guidance: local correctness can coexist with a mistaken view of the paper's purpose.

Only after the argument is coherent does line editing become efficient. Useful passes remove repeated material, replace vague claims with concrete comparisons, shorten descriptions of machinery, and improve the connection between paragraphs. They preserve numbers, units, uncertainty definitions, notation, and attribution. A proposed stylistic change that strengthens a scientific claim is a substantive revision and must be reviewed as such. Repeated generation without a clear diagnosis tends to produce more text instead of a better explanation.

The human author remains responsible for the final document, including passages the agent drafted. The goal is a paper whose claims, priorities, and reasoning the author can explain under questioning. In that sense, writing is also a test of scientific understanding: difficulty explaining the main point may reveal an unsettled interpretation or an unnecessary calculation. Effective assistance improves the argument while preserving this connection between the scientist and the work.

\section{Conclusions}\label{sec:conclusions}
We have described a living approach to reliable scientific work with agents, organized around numerical documentation, computational reliability, literature attribution, and scientific writing. Generated result macros and provenance records make numerical claims traceable and prevent defined forms of transcription drift. Tool execution, independent comparisons, and adversarial review supply evidence that a result survives relevant challenges. Citation tools and direct reading connect claims to genuine and appropriate previous work. These practices build on familiar standards of scientific work and review.

Human oversight remains essential to deciding which checks matter and what the results mean. Physical intuition and inspection of plots can expose failures outside an automated test suite, while an agent can turn those concerns into explicit calculations. Writing remains the least mechanical component: it requires context, priorities, capable models, sustained revision, and an author who owns the scientific argument. The guide and its supporting repository will evolve as those capabilities and practices improve.

\begin{acknowledgments}
This working manuscript and its public repository were assembled with OpenAI Codex under human direction, using previous human--agent research work and the explicit argument supplied for this revision. The agent retrieved policy and literature sources, reorganized and drafted text, and implemented and checked the accompanying synthetic tooling. All repository commits are performed by AI agents. This disclosure records the work performed; final human review and approval remain necessary before submission. Detailed model-version information and a final verification record will accompany a submitted version.
\end{acknowledgments}

\section*{Data and code availability}
The living guide, source, portable skills, and synthetic example are available in the \href{https://github.com/oshaughnessy-junior/whitepaper-transparent-agentic-science}{public companion repository}. The two scientific examples are the main papers in \href{https://github.com/oshaughnessy-junior/popsynth_hyperpipe_paper}{popsynth\_hyperpipe\_paper} and \href{https://github.com/oshaughnessy-junior/gwk-hierarchical-initial}{gwk-hierarchical-initial}; both repositories currently require permission to access. Their main preprints are publicly available~\cite{SyntheticUniverses,HierarchicalMethods}. The tooling source \href{https://github.com/oshaughnessy-junior/agent-skills}{agent-skills} also currently has restricted access. Public textual examples are not represented as new end-to-end replications of the scientific analyses.

\appendix
\section{An executable consistency example}\label{app:tooling}
A small synthetic example isolates the numerical documentation mechanism from the astrophysical applications. It evaluates
\begin{equation}
I=\int_{0}^{1}x^2\,\mathrm{d}x
\end{equation}
by composite midpoint quadrature and an antiderivative. Exact rational accumulation removes stochastic variation from the quadrature. The shared mathematical specification and runtime limit the independence of the comparison.

\begin{table}[tb]
\caption{Synthetic consistency example. The result record generates the table and numerical prose together; it records input and producer identities. OpenAI Codex implemented the generator and its tests under human direction.}
\label{tab:tooling}
\centering
\begin{tabular}{lr}
\toprule
Quantity & Value (dimensionless)\\
\midrule
Midpoint estimate & $\result{IntegralEstimate}$\\
Antiderivative reference & $\result{IntegralReference}$\\
Absolute discrepancy & $\result{IntegralError}$\\
Acceptance threshold & $\result{IntegralTolerance}$\\
\bottomrule
\end{tabular}

\end{table}
The estimate is $\result{IntegralEstimate}$, with absolute discrepancy $\result{IntegralError}$ below the prescribed tolerance $\result{IntegralTolerance}$ (Table~\ref{tab:tooling}). These quantities are inserted from the same generated record. Tests deliberately change inputs and producer identity, corrupt generated artifacts, introduce an undefined result or untracked source, and insert a typed numerical result. A wrong calculation must fail the numerical acceptance test, while refinement must improve the approximation in the expected way. This demonstrates specific checks, not a general proof of scientific correctness or a benchmark of agent performance.

\section{Living resources}\label{app:resources}
The public repository separates the scientific guide from its operational material. The latter includes the earlier agent-facing document, reusable skills for provenance, simulation validation, adversarial review, and scientific writing, and the synthetic example's regeneration and checking tools. These resources support the argument in the paper without requiring the reader to follow an agent's execution history. New lessons can be added with their evidence and limitations, while scientific claims remain tied to a particular reviewed version.

\begin{thebibliography}{99}
\bibitem{Hao2026}
Q.~Hao, F.~Xu, Y.~Li, and J.~Evans,
\href{https://doi.org/10.1038/s41586-025-09922-y}{Artificial intelligence tools expand scientists' impact but contract science's focus}, Nature \textbf{649}, 1237--1243 (2026).
\bibitem{Boiko2023}
D.~A.~Boiko, R.~MacKnight, B.~Kline, and G.~Gomes,
\href{https://doi.org/10.1038/s41586-023-06792-0}{Autonomous chemical research with large language models}, Nature \textbf{624}, 570--578 (2023).
\bibitem{APS2026}
American Physical Society, \href{https://journals.aps.org/authors/appropriate-use-ai-tools}{Appropriate Use of AI Tools}, Physical Review editorial policy (accessed October 23, 2026).
\bibitem{AAS2026}
American Astronomical Society, \href{https://journals.aas.org/author-llm-guidelines}{Author Guidelines for Use of AI and LLMs in Manuscript Preparation} (initially posted September 9, 2026; accessed October 23, 2026).
\bibitem{arXivAI}
arXiv, \href{https://info.arxiv.org/help/moderation/index.html#policy-for-authors-use-of-generative-ai-language-tools}{Policy for authors' use of generative AI language tools}, Content Moderation (accessed October 23, 2026).
\bibitem{SyntheticUniverses}
R.~O'Shaughnessy, R.~Mechum, M.~Qazalbash, Z.~Rosenberg, and M.~Zeeshan,
\href{https://arxiv.org/abs/2609.05996}{Physics-based phenomenological modeling of binary black hole hierarchical formation 1: Synthetic universes from globular cluster simulations for GWTC}, arXiv:2609.05996 (2026).
\bibitem{HierarchicalMethods}
R.~O'Shaughnessy, M.~Zeeshan, and M.~Qazalbash,
\href{https://arxiv.org/abs/2609.06728}{Physics-based phenomenological modeling of binary black hole hierarchical formation 2: Autodifferentiable functional inference of hierarchical compact-binary populations}, arXiv:2609.06728 (2026).
\bibitem{GWTCRelease}
Gravitational Wave Open Science Center, \href{https://doi.org/10.7935/b024-1886}{GWTC-3 Data Release Documentation}.
\end{thebibliography}
\end{document}
